\addcontentsline{toc}{subsubsection}{Solution 66. The Irrational Iteration Root}

\begin{dlsummary}[title={\bfseries Solution 66. The Irrational Iteration Root},breakable]

\begin{proof}
The unique solution is
\[
\boxed{f(x)=x^{\sqrt2}\qquad(x>0).}
\]

\paragraph{1. Recover strict monotonicity.}

If \(b>a>0\), superadditivity and positivity give
\[
f(b)=f(a+(b-a))
\ge f(a)+f(b-a)>f(a).
\]
Thus \(f\) is strictly increasing and, in particular, injective.

\paragraph{2. Force multiplicativity through composition.}

By submultiplicativity,
\[
f(xy)\le f(x)f(y).
\]
Applying the increasing function \(f\),
\[
f(f(xy))\le f(f(x)f(y)).
\]
Applying submultiplicativity again, now to \(f(x),f(y)\),
\[
f(f(x)f(y))
\le f(f(x))f(f(y)).
\]
The iteration identity therefore produces the squeeze
\[
\boxed{
x^2y^2
\le f(f(x)f(y))
\le x^2y^2.
}
\]
Hence
\[
f(f(x)f(y))=f(f(xy)).
\]
Injectivity gives
\[
\boxed{f(xy)=f(x)f(y).} \tag{1}
\]

This is the main breakthrough: the multiplicative inequality cannot remain strict anywhere.

\paragraph{3. Obtain the normalisation and rational powers.}

Taking \(x=y=1\) in (1), positivity yields
\[
f(1)=1.
\]
Also,
\[
f(2)\ge2f(1)=2.
\]
Define
\[
p=\log_2 f(2),
\]
so \(p\ge1\).

Multiplicativity gives, for every rational number \(r\),
\[
\boxed{f(2^r)=f(2)^r=2^{pr}.} \tag{2}
\]
For positive rational exponents, this follows by taking integer powers; reciprocal inputs handle negative exponents.

\paragraph{4. Determine every positive real input.}

Fix \(x>0\). Choose rational sequences \(r_n,s_n\) such that
\[
r_n<\log_2x<s_n,
\qquad r_n,s_n\longrightarrow\log_2x.
\]
Monotonicity and (2) give
\[
2^{pr_n}\le f(x)\le2^{ps_n}.
\]
The explicit bounds converge to \(x^p\). Therefore
\[
\boxed{f(x)=x^p\qquad(x>0).} \tag{3}
\]
No continuity of \(f\) was assumed.

Now the iteration identity becomes
\[
x^{p^2}=x^2\qquad(x>0).
\]
Thus \(p^2=2\), and \(p\ge1\) forces
\[
p=\sqrt2.
\]

\paragraph{5. Verify the solution.}

For \(f(x)=x^{\sqrt2}\), multiplicativity is immediate, and
\[
f(f(x))=(x^{\sqrt2})^{\sqrt2}=x^2.
\]

To verify superadditivity, note that for \(p\ge1\),
\[
\begin{aligned}
(x+y)^p
&=x(x+y)^{p-1}+y(x+y)^{p-1}\\
&\ge x\,x^{p-1}+y\,y^{p-1}\\
&=x^p+y^p.
\end{aligned}
\]
Taking \(p=\sqrt2\) completes the verification.
\end{proof}
\end{dlsummary}

\begin{takeaways}
\begin{itemize}
\item \textbf{Use inequalities at two levels:} First at \(x,y\), then at their function values.
\item \textbf{Iteration can close a squeeze:} The composition identity makes the two endpoints coincide exactly.
\item \textbf{Derive multiplicativity rather than guessing it:} This is the essential classification step.
\item \textbf{Separate regularity from algebra:} Positivity and superadditivity supply order; multiplicativity and rational approximation then identify the power.
\end{itemize}
\end{takeaways}
